% ============================================================
% analy 报告 — 规范模板 v2（xelatex + ctex）
% 主题: PyQUDA 软件库全面分析
% 编译: xelatex 两遍; 交付前 Overfull / Float too large 均为 0
% ============================================================
\documentclass[UTF8]{ctexart}
\usepackage[margin=2.5cm]{geometry}
\usepackage{graphicx}
\usepackage{amsmath}
\usepackage{microtype}
\usepackage{listings}
\usepackage{xcolor}
\usepackage{booktabs}
\usepackage{longtable}
\usepackage{xltabular}
\usepackage{tabularx}
\usepackage{hyperref}
\usepackage{fancyhdr}
\usepackage[most]{tcolorbox}
\usepackage{titlesec}

\definecolor{accent}{HTML}{1F4E79}
\definecolor{evidence}{HTML}{1E8449}
\definecolor{reference}{HTML}{8C6D1F}
\definecolor{conclusion}{HTML}{8B1A1A}
\definecolor{codebg}{HTML}{F4F6F7}
\definecolor{struct}{HTML}{3B3B98}
\definecolor{relation}{HTML}{0E7C7B}
\definecolor{idea}{HTML}{B03A2E}
\definecolor{warn}{HTML}{D35400}
\definecolor{hl}{HTML}{FDF6E3}

\setlength{\emergencystretch}{3em}
\hypersetup{colorlinks=true,linkcolor=accent,urlcolor=relation}

\titleformat{\section}{\Large\bfseries\color{accent}}{\thesection}{1em}{}
\titleformat{\subsection}{\large\bfseries\color{struct}}{\thesubsection}{1em}{}
\titleformat{\subsubsection}{\normalsize\bfseries\color{struct!70!black}}{\thesubsubsection}{1em}{}

\newtcolorbox{conclbox}{breakable,colback=conclusion!6,colframe=conclusion,title=结论}
\newtcolorbox{refbox}{breakable,colback=reference!6,colframe=reference,title=参考背景}
\newtcolorbox{ideabox}{breakable,colback=idea!6,colframe=idea,title=项目思路}
\newtcolorbox{warnbox}{breakable,colback=warn!6,colframe=warn,title=局限与风险}
\newtcolorbox{keybox}{breakable,colback=hl,colframe=accent,title=要点}

\lstset{
  basicstyle=\ttfamily\footnotesize,
  backgroundcolor=\color{codebg},
  frame=single,
  language=Python,
  firstnumber=auto,
  keywordstyle=\color{accent}\bfseries,
  commentstyle=\color{evidence!70!black},
  stringstyle=\color{struct},
  breaklines=true,
  breakatwhitespace=false,
  postbreak=\mbox{\textcolor{accent}{$\hookrightarrow$}\space},
  keepspaces=true,
  columns=flexible,
  numbers=left,numberstyle=\tiny\color{struct!60!black},
}

\title{PyQUDA 软件库全面分析报告}
\author{opencode analy}
\date{2026-08-17}

\begin{document}
\maketitle

\begin{abstract}
本报告对 QUDA 的 Python/Cython 包装库 \textbf{PyQUDA}（位于
\url{/root/PyQCU/refer/git-rep/PyQUDA}）进行只读全面分析。PyQUDA 以 Cython 封装 QUDA 的
C API，利用 CuPy/PyTorch 的 GPU 线性代数，覆盖 Wilson/Clover/HISQ 费米子作用量、multigrid、
HMC、规范场 smearing（APE/stout/HYP）、gradient flow、特征求解与夸克传播子计算，并支持多种
规范场 I/O 格式与多 GPU（mpi4py）。报告覆盖三视角：\textcolor{struct}{主体结构}
（核心桥/通信/工具/I-O/插件/示例分层）、\textcolor{relation}{各部分关系}
（`pyquda\_pyx.py 生成管线 → src/quda.pxd + quda.pyx → pyquda/dirac 封装` 链、
`pyquda\_utils → pyquda\_core` 高层依赖链）、\textcolor{idea}{项目思路}
（自动生成 Cython 桥 + 物理对象封装分离 + 多后端数组抽象）。并对"从源码安装到运行物理示例"
完整工作流按 A 框架 15 步重现。本快照非独立 git 仓库（共享 \texttt{/root/PyQCU/.git}）。
\end{abstract}

\tableofcontents

\section{仓库概览}

\subsection{目录结构}
\begin{lstlisting}[language=bash]
PyQUDA/
|-- README.md / pyproject.toml / setup.py     # 说明与打包（PyQUDA-Utils）
|-- pyquda_core/
|   |-- pyquda/          # 核心封装：field/gamma/hmc/enum_quda.in.py/quda.pyi
|   |   |-- action/      # GaugeAction/CloverWilsonAction/HISQAction
|   |   |-- dirac/       # WilsonDirac/CloverWilsonDirac/StaggeredDirac/HISQDirac/GaugeDirac/general
|   |   `-- src/         # quda.in.pxd / quda.in.pyx 模板
|   |-- pyquda_comm/     # MPI 网格/数组后端抽象/场类/pointer Cython
|   |-- quda/            # QUDA 头文件快照（quda.h/enum_quda.h 等 5 文件）
|   |-- pycparser/       # vendored C 解析器（子模块）
|   `-- pyquda_pyx.py    # Cython 桥自动生成器
|-- pyquda_io/           # 多格式 I/O（Chroma/MILC/ILDG/KYU/XQCD/NERSC/OpenQCD/NPY）
|-- pyquda_utils/        # 高层工具：core/convert/fft/phase/source/hmc_param/io 等
|-- pyquda_plugins/      # 插件框架：pycontract（夸克收缩）、pygwu（Overlap）
|-- examples/            # 淬火 HMC / 传播子 / 介子-质子 2pt / PCAC / 色散
|-- tests/               # pytest 套件
|-- docs/                # 本报告与历史 pure 报告
\end{lstlisting}

\subsection{规模统计与 git 信息}
\begin{table}[htbp]\centering
\begin{tabular}{ll}
\toprule
指标 & 实测值 \\
\midrule
文件总数 & 492（含 vendored pycparser 198） \\
非 vendored Python 源 & 108 文件 \\
pyquda\_core 合计 & 235 文件 / 57737 行 \\
\hspace{1em}pyquda/ & 25 文件 / 8184 行 \\
\hspace{1em}pyquda\_comm/ & 7 文件 / 2931 行 \\
\hspace{1em}quda/（头文件） & 5 文件 / 2747 行 \\
\hspace{1em}pycparser（vendored） & 198 文件 / 43875 行 \\
pyquda\_utils & 20 文件 / 3720 行 \\
pyquda\_io & 13 文件 / 1425 行 \\
pyquda\_plugins & 151 文件 / 12635 行 \\
examples & 8 文件 / 2077 行 \\
独立 git 仓库 & 否（共享 /root/PyQCU/.git，快照形态） \\
\bottomrule
\end{tabular}
\caption{仓库实测统计（数据来源: find/wc/git 实测）}
\end{table}

\section{主体结构}
\begin{xltabular}{\textwidth}{lXX}
\toprule
\textcolor{struct}{层次} & 目录 & \textcolor{struct}{职责} \\
\midrule
\endhead
桥接生成层 & \texttt{pyquda\_core/pyquda\_pyx.py} & 解析 QUDA 的 quda.h，生成 Cython 桥（quda.pxd/quda.pyx/enum\_quda）\\
Cython 桥层 & \texttt{pyquda\_core/pyquda/src/} & quda.in.pxd/quda.in.pyx 模板，经生成后为 src/quda.pxd/quda.pyx \\
物理封装层 & \texttt{pyquda\_core/pyquda/} & 场类型/$\gamma$ 矩阵/HMC/作用量/Dirac 算子高层 Python API \\
通信/数组层 & \texttt{pyquda\_core/pyquda\_comm/} & MPI 网格、numpy/cupy/dpnp/torch 后端抽象、场搬运、halo \\
头文件快照层 & \texttt{pyquda\_core/quda/} & QUDA 头文件（quda.h 1859 行、enum\_quda.h 643 行等）\\
解析依赖层 & \texttt{pyquda\_core/pycparser/} & vendored C 解析器（pycparser），桥生成使用 \\
高层工具层 & \texttt{pyquda\_utils/} & 求解封装（invert）、源构造、FFT、相位、gamma、HMC 参数、统一 I/O \\
I/O 层 & \texttt{pyquda\_io/} & 多格式规范场/传播子读写 \\
插件层 & \texttt{pyquda\_plugins/} & 夸克收缩（pycontract）、Overlap 算子（pygwu）、插件代码生成器 \\
示例层 & \texttt{examples/} & 淬火 HMC、夸克传播子、介子/质子 2pt、PCAC、色散关系 \\
测试层 & \texttt{tests/} & pytest 套件 \\
\bottomrule
\caption{主体结构分层（数据来源: 全库索引实测）}
\end{xltabular}

\begin{keybox}
\textbf{要点：}PyQUDA 的核心创新在于"自动生成 Cython 桥"——用 pycparser 解析 QUDA 头文件，
动态生成 Cython 声明与实现，避免手写巨量桥代码；物理 API 以 Python 类（\texttt{FermionDirac}
子类族）面向用户。
\end{keybox}

\section{各部分关系}

\subsection{桥生成链（核心机制）}
\begin{keybox}
\textbf{依赖主链:} \textcolor{relation}{\texttt{QUDA\_PATH/include/quda.h}
$\xrightarrow{\texttt{pyquda\_pyx.py:140 build\_pyquda\_pyx}}$
\texttt{src/quda.pxd + src/quda.pyx + enum\_quda.py + quda\_define.py}}
\end{keybox}
证据：\texttt{pyquda\_pyx.py:140}（入口）、\texttt{:186-207}（pycparser 解析收集 Quda 枚举与
Param 结构字段）、\texttt{:378-426}（按字段类型注入模板片段）、\texttt{:428-434}
（写出产物）；\texttt{setup.py:8-10} 在构建时调用。

\subsection{物理封装依赖链}
\begin{keybox}
\textbf{依赖主链:} \textcolor{relation}{\texttt{pyquda/dirac/wilson.py}
$\rightarrow$ \texttt{general.py}（newQudaGaugeParam/newQudaMultigridParam/newQudaInvertParam）
$\rightarrow$ \texttt{src/quda.pyx}（桥函数）$\rightarrow$ QUDA 库}
\end{keybox}
证据：\texttt{pyquda/dirac/wilson.py:19-47}（WilsonDirac 构造：$\kappa$ 计算、三组 Quda 参数、
multigrid 支持）；\texttt{dirac/general.py:283}（newQudaMultigridParam）、\texttt{:478}
（class Multigrid）。

\subsection{高层工具依赖链}
\begin{keybox}
\textbf{依赖主链:} \textcolor{relation}{\texttt{pyquda\_utils/core.py}（invert/getDirac 等）
$\rightarrow$ \texttt{pyquda\_core/pyquda/dirac/*} $\rightarrow$ \texttt{src/quda.pyx}}
\end{keybox}
证据：\texttt{pyquda\_utils/core.py:61}（invert）、\texttt{:430}（getWilson）、\texttt{:388}
（getDirac）；examples 统一经 \texttt{pyquda\_utils} 入口（\texttt{examples/1\_Quenched\_HMC.py:6-7}）。

\subsection{文件依赖关系汇总}
\begin{xltabular}{\textwidth}{lXX}
\toprule
\textcolor{relation}{源（A）} & 目标（B） & 关系类型 \\
\midrule
\endhead
\texttt{pyquda\_pyx.py} & \texttt{src/quda.pxd}, \texttt{src/quda.pyx} & 生成 \\
\texttt{src/quda.in.pxd/.pyx} & \texttt{src/quda.pxd/.pyx} & 模板→产物 \\
\texttt{quda.in.pxd}（cdef extern）& \texttt{quda.h} & 声明引用 \\
\texttt{pyquda/dirac/*.py} & \texttt{general.py}, \texttt{src/quda.pyx} & 调用 \\
\texttt{pyquda\_utils/core.py} & \texttt{pyquda\_core/*} & 依赖 \\
\texttt{pyquda\_io/*} & 格式 I/O & 实现 \\
\texttt{pyquda\_plugins/plugin\_pyx.py} & 插件 pyx & 生成 \\
\texttt{examples/*.py} & \texttt{pyquda\_utils} & 使用 \\
\texttt{setup.py} & \texttt{pyquda\_pyx.py} & 构建期调用 \\
\bottomrule
\caption{各部分关系汇总（数据来源: 源码实测）}
\end{xltabular}

\section{项目思路}

\subsection{设计意图}
PyQUDA 的设计动机：格点 QCD 研究者希望在 Python 生态中获得 QUDA 的性能，同时保留
CuPy/PyTorch 的便利与多 GPU（MPI）支持（\texttt{README.md:3-7}）。核心难点是 QUDA C API
庞大，故采用\textbf{自动生成 Cython 桥}避免手工维护；物理 API 层以面向对象封装
（\texttt{FermionDirac} 族）向用户呈现统一接口，多后端数组抽象（numpy/cupy/dpnp/torch，
\texttt{pyquda\_comm/array.py}）保证移植性。

\begin{ideabox}
\textbf{项目思路核心总结：}一个"自动桥 + 对象封装 + 多后端数组"的 QUDA Python 前端——
机器生成 Cython 桥保证与 QUDA API 同步，物理层提供科研友好接口，数组后端抽象换取
CPU/GPU/Intel GPU 可移植。
\end{ideabox}

\subsection{演进脉络}
\texttt{README.md:7} 声明基于 QUDA `develop` 分支、兼容 pull/1489 之后任意 commit；功能清单
（\texttt{README.md:9-42}）反映从基础 Wilson 反演逐步扩展到 multigrid、HMC、smearing、
gradient flow 的演进。\texttt{pyquda.pyi}（2145 行，手动维护）与
\texttt{README.md:100} 注明函数签名与 docstring 需人工更新，体现"自动生成 + 人工维护 stub"
的双轨机制。插件生态（pycontract、pygwu，\texttt{README.md:74-79}）为后续扩展路径。

\subsection{典型工作流}
一次物理计算的标准流程：
\begin{enumerate}
  \item 安装：\texttt{export QUDA\_PATH=...}；\texttt{pip install pyquda pyquda-utils} 或源码
        \texttt{setup.py build\_ext --inplace}（\texttt{README.md:44-64}）；
  \item \texttt{core.init()} 初始化 QUDA（\texttt{examples/1\_Quenched\_HMC.py:9}）；
  \item 建 \texttt{LatticeInfo}、载规范场、选 Dirac 算子（\texttt{getDirac}），反演/传播子
        计算（\texttt{\nolinkurl{pyquda_utils/core.py}}）；
  \item 物理量组装（2pt/PCAC/色散）与输出。
\end{enumerate}

\section{主题分析——工作全流程重现（A 代码工作框架）}

\subsection{A1 目的与前期准备}
\textbf{目的：}在 Python 中调用 QUDA 做格点 QCD 计算（反演、传播子、HMC、smearing）。
\textbf{前置条件：}CUDA 环境、安装好的 QUDA（\texttt{QUDA\_PATH} 指向 usqcd 安装，
\texttt{README.md:46-49}）、Python 包（CuPy、mpi4py、numpy）。构建需 pycparser 子模块
（\texttt{.gitmodules:1-3}）。

\subsection{A2 输入是什么}
\begin{itemize}
  \item QUDA 头文件（\texttt{QUDA\_PATH/include/quda.h}）——桥生成输入
        （\texttt{pyquda\_pyx.py:142-144}）；
  \item 物理参数：格子尺寸、质量/$\kappa$、容差、几何块大小（multigrid）、HMC 参数（$\beta$、loop 路径）。
\end{itemize}

\subsection{A3 输出应该是什么}
物理结果：夸克传播子、两点函数、PCAC 质量、色散关系、HMC 轨迹与接受率，以及 I/O 写出的
多格式文件。

\subsection{A4 整体框架}
构建：\texttt{setup.py} $\rightarrow$ \texttt{build\_pyquda\_pyx} 生成桥 $\rightarrow$
编译 3 个扩展（pointer/quda/malloc\_quda，\texttt{setup.py:25-53}）；
运行：\texttt{examples/*.py} $\rightarrow$ \texttt{pyquda\_utils} $\rightarrow$
\texttt{pyquda/dirac} $\rightarrow$ \texttt{src/quda.pyx} $\rightarrow$ QUDA。

\subsection{A5 关键细节}
\begin{itemize}
  \item \textbf{参数映射：}Python 对象 $\leftrightarrow$ QUDA Quda*Param 结构体字段
        （\texttt{quda.in.pyx} 中 cdef class QudaGaugeParam 等）；
  \item \textbf{multigrid：}\texttt{\nolinkurl{newQudaMultigridParam}} 的
        geo\_block\_size 规范化（\texttt{\nolinkurl{pyquda_utils/geo_block_size.py:9}}）
        与 Multigrid 对象生命周期（\texttt{\nolinkurl{dirac/general.py:478}}）；
  \item \textbf{基转换：}Dirac--Pauli $\leftrightarrow$ DeGrand--Rossi 基
        （\texttt{pyquda\_utils/io/\_\_init\_\_.py}）；
  \item \textbf{CRC 校验：}SciDAC/MILC/NERSC checksum（\texttt{pyquda\_utils/checksum.py:15}）。
\end{itemize}

\subsection{A6 任务如何正确执行}
\begin{enumerate}
  \item \texttt{git clone --recursive} 拉取 pycparser 子模块（\texttt{README.md:59}）；
  \item \texttt{cd pyquda\_core \&\& python3 setup.py build\_ext --inplace}（\texttt{README.md:87-89}）；
  \item \texttt{python examples/1\_Quenched\_HMC.py} 运行淬火 HMC。
\end{enumerate}

\subsection{A7 实际执行情况}
\textcolor{warn}{未验证}：本快照无 QUDA 安装与构建产物，无法实测桥生成与运行。

\subsection{A8 实际输出情况}
\textcolor{warn}{未验证}：无实测输出；预期为物理量数值与轨迹日志。

\subsection{A9 结果分析}
\textcolor{warn}{未验证}：无法实测。\texttt{examples/1\_Quenched\_HMC.py:12-26} 给出
$\beta=6.20$、$u_0=0.855453$、16$^3\times$32、2000 轨迹的参考参数，可作为未来验证基准。

\subsection{A10 综合分析}
\textbf{代码--对象映射：}本库代码符号对应物理对象的映射表：
\begin{xltabular}{\textwidth}{lXX}
\toprule
\textcolor{struct}{代码符号} & 物理/算法对象 & 含义 \\
\midrule
\endhead
\texttt{WilsonDirac} & Wilson 费米子算子 & $\kappa=1/[2(m+1+(N_d-1)/\xi)]$（\texttt{wilson.py:19}）\\
\texttt{CloverWilsonDirac} & Clover 改进算子 & 含 clover 项，支持 multigrid \\
\texttt{HISQDirac} & HISQ 作用量 & 高精度 staggered \\
\texttt{LatticeGauge} & $U_\mu(x)$ 规范场 & \texttt{pyquda\_comm/field.py} \\
\texttt{LatticeFermion} & $\psi(x)$ 费米子场 & 同上 \\
\texttt{GaugeAction} & 规范作用量 & plaquette 环路（\texttt{examples/1:14-24}）\\
\texttt{HMC} & HMC 驱动 & 轨迹积分 + 接受率（\texttt{hmc.py:262}）\\
\texttt{Multigrid} & 多重网格预条件子 & \texttt{dirac/general.py:478} \\
\texttt{mesonTwoPoint} & 介子两点函数 & \texttt{pycontract/__init__.py:19} \\
\texttt{checkGauge} & 西性/CRC 校验 & \texttt{checksum.py:15} \\
\bottomrule
\caption{代码符号 $\leftrightarrow$ 对象映射（数据来源: pyquda_core 与 examples 实测）}
\end{xltabular}

\subsection{A11 结尾总结}
\begin{conclbox}
PyQUDA 是"自动生成 Cython 桥 + 面向对象物理封装 + 多后端数组"的 QUDA Python 前端，
覆盖从反演到 HMC 的完整科研流程，是 PyQCU 生态中 Python 侧与 C++ QUDA 交互的重要参照。
\end{conclbox}

\subsection{A12 结尾评价}
\textbf{优点：}桥自动生成显著降低维护成本；物理 API 清晰统一；I/O 格式覆盖全；
多 GPU 支持。\textbf{可改进：}依赖 pycparser 子模块与精确 QUDA 版本；pyi stub 需人工维护；
插件依赖外部扩展库（contract/gwu）。

\subsection{A13 补充说明}
\textcolor{warn}{局限：}本快照无 QUDA 安装，A7/A8/A9 未实测；非独立 git 仓库。

\subsection{A14 参考源}
本节证据汇总见第 7 节参考源清单。

\subsection{A15 附录与附件（如有）}
\url{/root/PyQCU/refer/git-rep/PyQUDA/docs/pure_pyquda_core_20260815.tex} 为历史核心剖析报告；
\texttt{examples/} 为完整物理示例。

\section{关键代码/文档片段}

\begin{lstlisting}[language=Python]
class WilsonDirac(FermionDirac):
    def __init__(self, latt_info, mass, tol, maxiter, multigrid=None):
        kappa = 1 / (2 * (mass + 1 + (latt_info.Nd - 1) / latt_info.anisotropy))
        super().__init__(latt_info)
        self.newQudaGaugeParam()
        self.newQudaMultigridParam(multigrid, mass, kappa)
        self.newQudaInvertParam(mass, kappa, tol, maxiter)
        self.setPrecision()
        self.setReconstruct()
\end{lstlisting}
参考源：\texttt{pyquda\_core/pyquda/dirac/wilson.py:10-25}（WilsonDirac 构造）

\begin{lstlisting}[language=Python]
def build_pyquda_pyx(pyquda_root, quda_path):
    fake_libc_include = os.path.join(pyquda_root, "pycparser", "utils", "fake_libc_include")
    quda_include = os.path.join(quda_path, "include")
    assert os.path.exists(fake_libc_include), f"{fake_libc_include} not found"
    print(f"Building pyquda wrapper from {os.path.join(quda_include, 'quda.h')}")
\end{lstlisting}
参考源：\texttt{pyquda\_core/pyquda\_pyx.py:140-144}（桥生成入口）

\begin{lstlisting}[language=Python]
class O2Nf1Ng0V(Integrator):
    R"""https://doi.org/10.1016/S0010-4655(02)00754-3
    BAB: Eq.(23), \xi=0"""
    def integrate(self, updateGauge, updateMom, t: float):
        dt = t / self.n_steps
        for _ in range(self.n_steps):
            updateMom(dt / 2)
            updateGauge(dt)
            updateMom(dt / 2)
\end{lstlisting}
参考源：\texttt{examples/1\_Quenched\_HMC.py:29-38}（O2 积分器）

\section{参考源清单}
\begin{xltabular}{\textwidth}{lXX}
\toprule
\textcolor{evidence}{参考源} & 类型 & 内容/作用 \\
\midrule
\endhead
\texttt{\nolinkurl{README.md:3-7}} & 仓库 & 项目定位与 CuPy/PyTorch 动机 \\
\texttt{\nolinkurl{README.md:9-42}} & 仓库 & 功能特性清单 \\
\texttt{\nolinkurl{README.md:44-64}} & 仓库 & 安装步骤（PyPI/源码）\\
\texttt{\nolinkurl{README.md:81-94}} & 仓库 & 开发模式（inplace）\\
\texttt{\nolinkurl{README.md:100}} & 仓库 & 维护说明（pyi 人工更新）\\
\texttt{\nolinkurl{pyquda_core/pyquda_pyx.py:140}} & 仓库 & 桥生成入口 \\
\texttt{\nolinkurl{pyquda_core/pyquda_pyx.py:186-207}} & 仓库 & pycparser 解析 quda.h \\
\texttt{\nolinkurl{pyquda_core/pyquda_pyx.py:378-426}} & 仓库 & 字段模板注入 \\
\texttt{\nolinkurl{pyquda_core/pyquda_pyx.py:428-434}} & 仓库 & 产物写出 \\
\texttt{\nolinkurl{pyquda_core/pyquda/dirac/wilson.py:10-25}} & 仓库 & WilsonDirac 构造 \\
\texttt{\nolinkurl{pyquda_core/pyquda/dirac/general.py:283}} & 仓库 & newQudaMultigridParam \\
\texttt{\nolinkurl{pyquda_core/pyquda/dirac/general.py:478}} & 仓库 & class Multigrid \\
\texttt{\nolinkurl{pyquda_core/pyquda/hmc.py:262}} & 仓库 & class HMC \\
\texttt{\nolinkurl{pyquda_core/pyquda/field.py:224-299}} & 仓库 & APE/stout/HYP/flow 高层方法 \\
\texttt{\nolinkurl{pyquda_core/pyquda_comm/array.py}} & 仓库 & 多后端数组抽象 \\
\texttt{\nolinkurl{pyquda_core/pyquda_comm/field.py}} & 仓库 & 场类与 halo 交换 \\
\texttt{\nolinkurl{pyquda_core/setup.py:8-10}} & 仓库 & 构建期调用桥生成 \\
\texttt{\nolinkurl{pyquda_core/setup.py:25-53}} & 仓库 & 三个扩展编译 \\
\texttt{\nolinkurl{pyquda_utils/core.py:61}} & 仓库 & invert 封装 \\
\texttt{\nolinkurl{pyquda_utils/core.py:430}} & 仓库 & getWilson 工厂 \\
\texttt{\nolinkurl{pyquda_utils/checksum.py:15}} & 仓库 & checksumSciDAC \\
\texttt{\nolinkurl{pyquda_utils/geo_block_size.py:9}} & 仓库 & multigrid 块大小规范化 \\
\texttt{\nolinkurl{pyquda_io/__init__.py:1-16}} & 仓库 & I/O 格式门面 \\
\texttt{\nolinkurl{pyquda_plugins/plugin_pyx.py:128}} & 仓库 & 插件头解析生成 \\
\texttt{\nolinkurl{pyquda_plugins/pycontract/__init__.py:19}} & 仓库 & mesonTwoPoint \\
\texttt{\nolinkurl{pyquda_plugins/pygwu/__init__.py:187}} & 仓库 & Overlap 算子 \\
\texttt{\nolinkurl{examples/1_Quenched_HMC.py:9-26}} & 仓库 & 淬火 HMC 参数 \\
\texttt{\nolinkurl{examples/1_Quenched_HMC.py:29-38}} & 仓库 & O2 积分器 \\
\texttt{\nolinkurl{.gitmodules:1-3}} & 仓库 & pycparser 子模块声明 \\
\bottomrule
\end{xltabular}

\section{结论与局限}
\begin{conclbox}
三视角结论：\textcolor{struct}{结构}——桥生成/物理封装/通信/工具/I-O/插件六层分明；
\textcolor{relation}{关系}——`pyquda\_pyx.py → src/quda.pyx → dirac/* → pyquda\_utils` 主链贯通，
自动生成与人工 stub 双轨；\textcolor{idea}{思路}——以自动生成化解 QUDA API 规模，以对象封装
提供科研接口，以多后端抽象保证可移植。
\end{conclbox}
\begin{warnbox}
\textcolor{warn}{未验证}：A7/A8/A9 实测执行步骤无法在本快照验证（无 QUDA 安装）；行号均经
read/grep 核实；本库非独立 git 仓库。
\end{warnbox}

\end{document}